\begin{onehalfspace}

Questa appendice raccoglie in forma estesa il contesto richiamato sinteticamente nel Capitolo~\ref{chap:introduzione}: i principi fisici del calcolo quantistico e la loro implementazione hardware, la rivoluzione algoritmica degli anni Novanta, il protocollo RSA e la sua vulnerabilità all'algoritmo di Shor, la crittografia post-quantistica e lo stato dell'arte dell'hardware quantistico.

\subsection{Dal Bit al Qubit: i Principi Fisici del Calcolo Quantistico}
\label{app:intro:qubit}

Per comprendere appieno le potenzialità e i limiti del calcolo quantistico, è necessario esaminare le differenze fisiche fondamentali tra l'unità di informazione classica e quella quantistica.

\subsubsection{Il Qubit e la Sovrapposizione}

Un computer classico elabora l'informazione attraverso \textbf{bit} -- unità fisiche che possono trovarsi in uno di due stati discreti, convenzionalmente 0 e 1, realizzati come livelli di tensione in circuiti elettronici. L'unità fondamentale del calcolo quantistico è invece il \textbf{qubit} (\textit{quantum bit}): un sistema fisico a due livelli che obbedisce alle leggi della meccanica quantistica.

Matematicamente, lo stato di un qubit è descritto da un vettore nello spazio di Hilbert $\mathcal{H} \cong \mathbb{C}^2$:
\begin{equation}
    \ket{\psi} = \alpha\ket{0} + \beta\ket{1}, \qquad \alpha,\beta \in \mathbb{C}, \quad |\alpha|^2 + |\beta|^2 = 1
    \label{eq:qubit_state}
\end{equation}
dove $\ket{0}$ e $\ket{1}$ sono i due stati della base computazionale, analoghi ai valori 0 e 1 del bit classico. La differenza fondamentale è che $\alpha$ e $\beta$ possono essere simultaneamente diversi da zero: il qubit si trova in una \textbf{sovrapposizione} dei due stati base, con probabilità $|\alpha|^2$ di essere misurato in $\ket{0}$ e probabilità $|\beta|^2$ di essere misurato in $\ket{1}$.

È cruciale sottolineare che la sovrapposizione non equivale a ignoranza sullo stato del sistema: il qubit non è ``segretamente'' in uno stato definito che non conosciamo. Si tratta di una proprietà fisica genuina, verificata sperimentalmente con precisione straordinaria, che non ha analogo nella fisica classica.

\subsubsection{La Sfera di Bloch}

Poiché la norma è unitaria e la fase globale è fisicamente irrilevante, lo stato di un qubit puro può essere parametrizzato da soli due angoli reali:
\begin{equation}
    \ket{\psi} = \cos\!\frac{\theta}{2}\ket{0} + e^{i\phi}\sin\!\frac{\theta}{2}\ket{1}, \qquad \theta \in [0,\pi],\quad \phi \in [0,2\pi)
    \label{eq:bloch_parametrization}
\end{equation}
Questo definisce un punto sulla superficie di una sfera unitaria -- la \textbf{sfera di Bloch} -- dove il polo nord corrisponde a $\ket{0}$, il polo sud a $\ket{1}$, e i punti equatoriali a sovrapposizioni equiprobabili con fasi diverse. Questa rappresentazione geometrica è di grande utilità pratica: le operazioni su un singolo qubit corrispondono a \textbf{rotazioni rigide} della sfera di Bloch, e il rumore -- come si vedrà nel Capitolo~\ref{chap:rumore} -- corrisponde alla contrazione progressiva del vettore di Bloch verso il centro, fino a ridurlo a zero (stato completamente misto, massima incertezza).

\subsubsection{L'Entanglement}

Quando due o più qubit interagiscono tramite porte quantistiche opportune, possono formare stati \textbf{entangled}: stati del sistema composto che non possono essere scritti come prodotto tensoriale di stati individuali. Il prototipo è lo stato di Bell:
\begin{equation}
    \ket{\Phi^+} = \frac{1}{\sqrt{2}}\bigl(\ket{00} + \ket{11}\bigr)
    \label{eq:bell_state}
\end{equation}
In questo stato, misurare il primo qubit in $\ket{0}$ collassa istantaneamente il secondo in $\ket{0}$, e viceversa per $\ket{1}$, indipendentemente dalla distanza fisica tra i due qubit. L'entanglement crea correlazioni impossibili da replicare con qualunque modello classico a variabili nascoste locali, come dimostrato teoricamente da Bell nel 1964 \cite{bell1964} e verificato sperimentalmente con crescente precisione negli anni successivi \cite{aspect1982}. Nel contesto del calcolo quantistico, l'entanglement è la risorsa che permette agli algoritmi di elaborare esponenzialmente più informazione rispetto ai circuiti classici: un registro di $n$ qubit entangled descrive uno stato nello spazio $\mathbb{C}^{2^n}$, inaccessibile a qualunque rappresentazione classica di dimensione polinomiale in $n$.

\subsubsection{L'Interferenza Quantistica}

Il terzo pilastro del calcolo quantistico è l'\textbf{interferenza}: la capacità di far sommare costruttivamente o annullare distruttivamente le ampiezze di probabilità di percorsi computazionali diversi. Le ampiezze $\alpha$ e $\beta$ in \eqref{eq:qubit_state} sono numeri complessi, e possono quindi avere segni opposti o sfasamenti arbitrari: quando si sommano, possono amplificarsi o cancellarsi esattamente come onde fisiche.

Un algoritmo quantistico ben progettato sfrutta sistematicamente questo meccanismo: attraverso una sequenza opportuna di porte quantistiche, si costruisce un'interferenza \textit{costruttiva} sui percorsi che portano alla risposta corretta e \textit{distruttiva} su quelli errati, concentrando la probabilità di misura sull'output desiderato. È esattamente questo il meccanismo alla base della ricerca quadraticamente accelerata di Grover e della fattorizzazione polinomiale di Shor.

\subsubsection{Implementazione Fisica: i Qubit Superconduttori di IBM}

Un qubit può essere realizzato con qualunque sistema fisico a due livelli: la polarizzazione di un fotone, lo spin di un elettrone, i livelli energetici di un atomo intrappolato. I \textbf{qubit superconduttori} utilizzati da IBM -- la piattaforma adottata nel presente lavoro -- sono circuiti elettrici macroscopici che operano nel regime quantistico grazie al raffreddamento a temperature criogeniche estreme.

Il componente chiave è la \textbf{giunzione di Josephson} \cite{josephson1962}: una sottile barriera isolante (spessa $\sim 1$--$2\,$nm) interposta tra due materiali superconduttori. Questo elemento introduce una non-linearità nell'energia del circuito che rompe la degenerazione dei livelli energetici, isolando i due stati computazionali $\ket{0}$ e $\ket{1}$ da tutti gli altri. Il risultato è il \textbf{transmon} \cite{koch2007} -- il tipo di qubit superconduttore adottato da IBM -- un oscillatore anelastico con frequenza tipica $4$--$6\,$GHz, controllato da impulsi a microonde.

Per operare nel regime quantistico, questi circuiti devono essere raffreddati a temperature dell'ordine di $15\,\text{mK}$ ($-273.13\,°$C), più fredde dello spazio cosmico a $2.7\,\text{K}$, tramite refrigeratori a diluizione. A queste temperature la resistenza elettrica scompare per effetto della superconduttività, eliminando la dissipazione termica che causerebbe la distruzione dello stato quantistico. 
\clearpage
I parametri operativi tipici dei processori IBM (serie Eagle, Heron) sono:

\begin{itemize}
    \item \textbf{Temperatura operativa}: $\sim 15\,\text{mK}$
    \item \textbf{Tempi di coerenza}: $T_1, T_2 \sim 50$--$500\,\mu$s
    \item \textbf{Durata gate a singolo qubit}: $\sim 50\,\text{ns}$
    \item \textbf{Durata gate a due qubit (ECR/CNOT)}: $\sim 300$--$500\,\text{ns}$
    \item \textbf{Fedeltà gate a singolo qubit}: $\sim 99.9$--$99.97\%$
    \item \textbf{Fedeltà gate a due qubit}: $\sim 99$--$99.7\%$
\end{itemize}

Il rapporto tra il tempo di coerenza e la durata di un gate determina quante operazioni possono essere eseguite prima che il qubit perda il proprio stato quantistico -- ed è proprio questo rapporto il principale collo di bottiglia per l'esecuzione di algoritmi profondi come Shor sui processori attuali.

% -----------------------------------------------
\subsection{La Rivoluzione Algoritmica degli Anni Novanta}
\label{app:intro:algoritmi}

Negli anni successivi alla proposta di Deutsch, la comunità scientifica si interrogò sulla effettiva esistenza di problemi che un computer quantistico potesse risolvere in modo \textit{probabilmente} più efficiente di qualunque macchina classica. La risposta arrivò in forma dirompente nel corso di pochi anni.

Nel 1992, Daniel Simon\footnote{L'\textbf{algoritmo di Simon} (1994, pubblicato nel 1997) fu il primo a dimostrare uno speedup \textit{esponenziale} rispetto a qualunque algoritmo classico probabilistico, stabilendo la prova di principio che i computer quantistici potessero essere genuinamente più potenti. Shor riconobbe che la sua tecnica si ispirava direttamente alle idee di Simon.} identificò un problema artificiale per il quale un algoritmo quantistico era esponenzialmente più veloce di qualunque algoritmo classico. Questo risultato rimase per alcuni anni una curiosità teorica, fino a quando, nel \textbf{1994}, Peter Shor pubblicò un risultato che scosse le fondamenta della sicurezza informatica mondiale: un algoritmo quantistico in grado di fattorizzare un intero di $n$ bit in tempo \textit{polinomiale} $O(n^3)$ \cite{shor1994}, a fronte della complessità sub-esponenziale ma super-polinomiale dei migliori algoritmi classici noti. Le implicazioni furono immediate e devastanti: l'intero edificio della crittografia a chiave pubblica -- \ac{RSA}, Diffie-Hellman, le firme digitali \ac{ECDSA} -- si fondava proprio sull'assunzione che la fattorizzazione fosse computazionalmente intrattabile. Un computer quantistico sufficientemente grande avrebbe reso obsoleta, in un solo colpo, l'intera infrastruttura crittografica su cui poggia Internet.

Due anni dopo, nel \textbf{1996}, Lov Grover presentò un algoritmo di tutt'altra natura, ma di eguale importanza \cite{grover1996}. Mentre Shor aveva risolto un problema specifico della teoria dei numeri, Grover affrontò una sfida universale: la \textbf{ricerca non strutturata}. Dato un insieme di $N$ elementi senza alcuna struttura o ordinamento, trovare quello che soddisfa una certa proprietà richiede, nel caso classico, un numero di interrogazioni proporzionale a $N$. Grover dimostrò che un algoritmo quantistico può fare lo stesso lavoro in sole $O(\sqrt{N})$ interrogazioni -- uno \textit{speedup quadratico} -- e che questo risultato è \textit{ottimale}: nessun algoritmo quantistico può fare di meglio \cite{zalka1999}. Sebbene meno spettacolare dello speedup esponenziale di Shor, il vantaggio di Grover è universale, applicabile a qualunque problema di ricerca, e porta con sé implicazioni profonde per la crittografia simmetrica, l'ottimizzazione combinatoria e la ricerca in grandi basi di dati.

Questi due algoritmi -- Shor e Grover -- rappresentano ancora oggi i pilastri dell'informatica algoritmica quantistica: il primo come dimostrazione del potere \textit{esponenziale} del calcolo quantistico su problemi strutturati, il secondo come strumento \textit{universale} di accelerazione quadratica applicabile al problema computazionale più fondamentale.

% -----------------------------------------------
\subsection{RSA e le Implicazioni Crittografiche dell'Algoritmo di Shor}
\label{app:intro:rsa}

Per comprendere appieno la portata rivoluzionaria dell'algoritmo di Shor è necessario esaminare la struttura matematica del sistema crittografico che esso minaccia. Il sistema \textbf{RSA}, proposto nel 1977 da Rivest, Shamir e Adleman \cite{rsa1978}, è la pietra angolare della crittografia a chiave pubblica e il fondamento di praticamente ogni protocollo sicuro su Internet.

\subsubsection{Struttura del Protocollo RSA}

La sicurezza di RSA si basa sulla \textbf{asimmetria computazionale} tra due operazioni inverse: moltiplicare due numeri primi grandi $p$ e $q$ per ottenere $N = p \cdot q$ è un'operazione elementare (complessità $O(n^2)$ per numeri a $n$ bit), mentre fattorizzare $N$ per recuperare $p$ e $q$ è computazionalmente intrattabile con gli algoritmi classici noti. La generazione delle chiavi segue i passi seguenti:

\begin{enumerate}
    \item Scegliere due primi grandi $p$ e $q$ (tipicamente a 1024--2048 bit ciascuno) e calcolare $N = p \cdot q$;
    \item Calcolare la funzione di Eulero $\phi(N) = (p-1)(q-1)$;
    \item Scegliere $e$ coprimo con $\phi(N)$ (convenzionalmente $e = 65537 = 2^{16}+1$, primo di Fermat);
    \item Calcolare $d \equiv e^{-1} \pmod{\phi(N)}$ tramite l'algoritmo di Euclide esteso.
\end{enumerate}

La \textbf{chiave pubblica} è la coppia $(N,\, e)$; la \textbf{chiave privata} è $d$. Cifrare un messaggio $m < N$ produce il crittogramma $c = m^e \bmod N$; decifrarlo richiede $m = c^d \bmod N$. La correttezza si fonda sul \textbf{teorema di Eulero} --- la generalizzazione del piccolo teorema di Fermat a moduli composti, secondo cui $m^{\phi(N)} \equiv 1 \pmod{N}$ per $m$ coprimo con $N$: essendo $ed \equiv 1 \pmod{\phi(N)}$, si ha $ed = 1 + k\phi(N)$ e quindi $(m^e)^d = m^{1+k\phi(N)} \equiv m \pmod{N}$. Tutta la sicurezza risiede nell'impossibilità di calcolare $d$ da $(N, e)$ senza conoscere la fattorizzazione di $N$.

\subsubsection{Perché l'Algoritmo di Shor Rompe RSA}

Il miglior algoritmo classico per la fattorizzazione di un intero a $n$ bit è il \textit{General Number Field Sieve} (GNFS) \cite{lenstra1993gnfs}, con complessità sub-esponenziale:
\begin{equation}
    T_{\text{GNFS}}(N) = O\!\left(\exp\!\left[c\cdot(\log N)^{1/3}(\log\log N)^{2/3}\right]\right)
    \label{eq:gnfs_complexity_intro}
\end{equation}
Per $N$ a 2048 bit, questa espressione implica circa $2^{112}$ operazioni classiche -- un numero astronomico, ben oltre la capacità di qualunque infrastruttura computazionale esistente o prevedibile in ambito classico. I migliori risultati pratici al 2020 hanno raggiunto la fattorizzazione di RSA-250 (829 bit) \cite{rsa250_2020}, impiegando circa 2700 anni-CPU di calcolo distribuito.

L'algoritmo di Shor riduce questa complessità a $O(n^3)$ -- un salto da sub-esponenziale a \textbf{polinomiale} -- rendendo la fattorizzazione di chiavi RSA-2048 un problema risolvibile in ore su un computer quantistico fault-tolerant con qualche migliaio di qubit logici. La struttura dell'algoritmo, descritta in dettaglio nel Capitolo~\ref{chap:shor}, sfrutta la Quantum Fourier Transform per trovare il periodo di una funzione modulare, problema che si riduce efficientemente alla fattorizzazione.

\clearpage
\subsubsection{Implicazioni Pratiche per la Sicurezza Informatica}

Le conseguenze sono sistemiche. RSA-2048 e le sue varianti sono il fondamento di:

\begin{itemize}
    \item \textbf{HTTPS e TLS}: proteggono oltre il 95\% del traffico web mondiale, incluse operazioni bancarie, sanitarie e governative;
    \item \textbf{Certificati X.509}: l'infrastruttura a chiave pubblica (PKI) su cui si basa l'autenticità di ogni sito web;
    \item \textbf{SSH}: il protocollo standard per l'accesso remoto sicuro ai server;
    \item \textbf{Firme digitali}: documenti legali, transazioni finanziarie, aggiornamenti software firmati.
\end{itemize}

Un computer quantistico scalabile renderebbe retroattivamente vulnerabili anche le comunicazioni \textit{storicamente registrate}, attraverso il cosiddetto attacco \textit{``harvest now, decrypt later''} (HNDL): avversari statali raccolgono oggi grandi volumi di traffico cifrato con RSA, conservandoli in attesa di poterli decifrare quando l'hardware quantistico raggiungerà la scala necessaria. Questo scenario non è teorico: rapporti di intelligence statunitensi e cinesi documentano campagne di raccolta massiva di traffico cifrato attive da almeno un decennio. Stime dell'NSA e del \ac{NIST} indicano che chiavi RSA-2048 potrebbero diventare vulnerabili entro il 2030--2035 con il proseguire dell'avanzamento tecnologico quantistico \cite{nsa_cnsa2022}.

\clearpage
\subsection{Crittografia Post-Quantistica: la Risposta dell'Industria}
\label{app:intro:pqc}

La consapevolezza della minaccia quantistica ha spinto il \ac{NIST} a lanciare nel 2016 un processo di standardizzazione internazionale per la \textbf{crittografia post-quantistica} (\ac{PQC}) -- algoritmi crittografici resistenti agli attacchi di computer quantistici, progettati per essere eseguiti su hardware classico convenzionale e quindi immediatamente distribuibili.

\subsubsection{Gli Standard NIST 2024}

Dopo otto anni di valutazione pubblica, crittoanalisi internazionale e tre successivi turni di selezione tra 69 candidati iniziali, nell'agosto 2024 il NIST ha pubblicato i primi tre standard definitivi \cite{nist2024pqc}:

\begin{itemize}
    \item \textbf{ML-KEM} (FIPS 203, ex CRYSTALS-Kyber): meccanismo di incapsulamento di chiavi basato sul problema \textit{Module Learning With Errors} (MLWE) su reticoli. È il sostituto designato per lo scambio di chiavi RSA e Diffie-Hellman nei protocolli TLS;

    \item \textbf{ML-DSA} (FIPS 204, ex CRYSTALS-Dilithium): schema di firma digitale basato sullo stesso fondamento matematico di ML-KEM. Sostituisce RSA e ECDSA per le firme digitali;

    \item \textbf{SLH-DSA} (FIPS 205, ex SPHINCS+): schema di firma basato esclusivamente su funzioni hash, senza assunzioni algebriche. Più conservativo per sicurezza, adottato come backup contro potenziali attacchi futuri ai reticoli.
\end{itemize}

La sicurezza di questi schemi si fonda su problemi matematici -- trovare vettori corti in reticoli ad alta dimensione, o invertire funzioni hash -- per i quali non sono noti algoritmi quantistici con speedup esponenziale. In particolare, l'algoritmo di Grover fornisce al più uno speedup quadratico sui problemi basati su hash, motivando l'aumento dei parametri di sicurezza (es. passaggio da 128 a 256 bit di sicurezza classica equivalente).

\subsubsection{La Sfida della Migrazione}

La transizione verso la crittografia post-quantistica è un'operazione di ingegneria su scala globale, paragonabile per complessità alla migrazione da IPv4 a IPv6, ma con urgenza maggiore. Le principali sfide includono:

\begin{itemize}
    \item \textbf{Eterogeneità dei sistemi}: miliardi di dispositivi embedded, sensori industriali e infrastrutture legacy con cicli di vita decennali che non riceveranno aggiornamenti software;
    \item \textbf{Overhead di banda}: il materiale crittografico post-quantistico è sensibilmente più voluminoso. Al livello di sicurezza più basso (ML-KEM-512), la chiave di incapsulamento occupa 800 byte e il crittogramma 768 byte, contro i circa 256 byte di una chiave pubblica RSA-2048: un fattore di crescita che pesa sui protocolli a bassa latenza e sui dispositivi con banda limitata;
    \item \textbf{Migrazione a doppio binario}: nella fase transitoria, molte implementazioni adottano schemi \textit{ibridi} (classico + post-quantistico in parallelo) per mantenere compatibilità con sistemi non aggiornati.
\end{itemize}

Aziende come Google (Chrome 131), Apple (iMessage PQ3) \cite{apple_pq3_2024} e Cloudflare hanno già avviato implementazioni ibride nei loro protocolli di produzione. La coesistenza di questa transizione con la ricerca attiva su algoritmi quantistici come Shor definisce il contesto strategico in cui si inserisce il presente lavoro: comprendere i limiti reali di Shor su hardware NISQ è rilevante non solo teoricamente, ma anche per calibrare con maggiore precisione le tempistiche e le priorità della transizione crittografica globale.

\clearpage
\subsection{Lo Stato dell'Arte: dall'Era NISQ ai Computer Fault-Tolerant}
\label{app:intro:nisq}

Nonostante la solidità teorica di questi risultati, la realizzazione pratica di computer quantistici di scala sufficiente ad eseguirli su istanze realmente significative incontra ostacoli formidabili. I sistemi quantistici fisici sono intrinsecamente fragili: qualunque interazione non controllata con l'ambiente esterno -- vibrazioni termiche, campi elettromagnetici, imperfezioni nei materiali -- provoca fenomeni di \textbf{decoerenza} -- ossia la perdita progressiva della coerenza quantistica causata dall'entanglement indesiderato tra il qubit e i gradi di libertà ambientali, che sopprime le interferenze quantistiche riducendo il sistema a una miscela statistica classica -- degradando lo stato quantistico in modo irreversibile prima che il calcolo possa completarsi. A differenza di un bit classico, che può essere protetto dalla ridondanza, un qubit non può essere copiato (teorema del no-cloning \cite{wootters1982}) né letto senza distruggerlo: questo rende la correzione degli errori quantistici un problema radicalmente diverso e molto più complesso di quella classica.

Il panorama tecnologico attuale è dominato da un insieme variegato di piattaforme hardware in competizione, ciascuna con vantaggi e svantaggi distinti:

\begin{itemize}
    \item \textbf{Qubit superconduttori} (IBM, Google, Rigetti): alta velocità dei gate ($\sim 50$--$500\,$ns), scalabilità su chip in silicio, ma tempi di coerenza limitati a $\sim 100$--$500\,\mu$s e necessità di raffreddamento a $15\,\text{mK}$;
    \item \textbf{Ioni intrappolati} (IonQ, Quantinuum): tempi di coerenza eccellenti ($>1\,$s), fedeltà dei gate superiore al $99.9\%$ anche su due qubit, ma gate lenti ($\sim 10$--$100\,\mu$s) e scalabilità più complessa;
    \item \textbf{Fotoni} (PsiQuantum, Xanadu): operano a temperatura ambiente, naturalmente adatti alla comunicazione quantistica, ma la generazione deterministica di fotoni entangled rimane una sfida aperta;
    \item \textbf{Spin in silicio} (Intel, UNSW): compatibili con i processi di fabbricazione CMOS esistenti, potenzialmente scalabili a milioni di qubit, ma con tecnologia di controllo ancora in fase di sviluppo.
\end{itemize}

Un momento simbolico di questa corsa fu raggiunto nel \textbf{2019}, quando il team di Google rivendicò il raggiungimento della cosiddetta \textit{quantum supremacy}\footnote{Il termine \textit{quantum supremacy}, coniato da John Preskill nel 2012, indica il punto in cui un computer quantistico esegue un calcolo specifico in un tempo che nessun supercomputer classico potrebbe eguagliare in modo pratico. La rivendicazione di Google con il processore Sycamore è rimasta controversa: IBM sostenne che il calcolo avrebbe potuto essere eseguito classicamente in 2.5 giorni anziché 10.000 anni, con un algoritmo migliorato. Questo dibattito evidenzia la difficoltà di definire e misurare il vantaggio quantistico.} con il processore Sycamore a 53 qubit \cite{arute2019}: un calcolo specifico fu completato in 200 secondi, mentre le stime attribuivano al più potente supercomputer classico un tempo di circa 10.000 anni per lo stesso compito. Nel 2023, IBM ha presentato il processore \textit{Condor} a 1.121 qubit e ha delineato una roadmap \cite{ibm_roadmap2025} verso sistemi con architettura fault-tolerant entro la fine del decennio, con il target \textit{Quantum Starling} previsto per il 2029.

Tuttavia, la distanza tra i dispositivi attuali e un computer quantistico \textit{fault-tolerant} -- in grado di eseguire Shor su istanze RSA reali -- è ancora enorme. I processori odierni operano nel regime cosiddetto \ac{NISQ}, un termine coniato da John Preskill nel 2018 \cite{preskill2018} per descrivere macchine con decine o centinaia di qubit fisici, affette da errori significativi che limitano la profondità dei circuiti eseguibili. In questo scenario, la ricerca si concentra su due direzioni parallele e complementari:

\begin{itemize}
    \item \textbf{\ac{QEC}}: tecniche per codificare qubit logici in qubit fisici ridondanti, in modo da rilevare e correggere gli errori durante il calcolo, a costo di un overhead fisico significativo (centinaia o migliaia di qubit fisici per qubit logico);

    \item \textbf{Noise mitigation}: tecniche che, senza richiedere overhead hardware, riducono l'effetto del rumore a livello di post-processing classico, estendendo la portata dei circuiti eseguibili su hardware NISQ.
\end{itemize}

\end{onehalfspace}
